\section{Related Work}\label{sec:related}

\paragraph{Learning with noisy labels.}
Surveys group methods for deep networks \citep{song2023learning} and for classical classifiers \citep{frenay2014classification} into a few families. Loss correction estimates a matrix of class-corruption probabilities and adjusts the loss \citep{patrini2017making}. Noise-tolerant losses are chosen so that the minimizer survives label noise \citep{ghosh2017robust}. Noise models can also sit inside logistic regression \citep{bootkrajang2014learning,lee2021individual}. The fourth family, sample selection, picks the training samples to trust. Our out-of-fold filter changes neither the loss nor the model and needs no noise model.

\paragraph{Sample selection in deep networks.}
In co-teaching, two networks exchange their small-loss samples \citep{han2018coteaching}. Co-teaching+ selects among the samples on which the two networks disagree \citep{yu2019how}. Area Under the Margin ranks samples by their training dynamics \citep{pleiss2020identifying}. These methods rely on deep-network training dynamics. We use the out-of-fold predictions of a linear model and do not compare against them.

\paragraph{Cross-validated filtering and confident learning.}
\citet{brodley1999identifying} filter a training set by removing the instances that classifiers misclassify, in single-algorithm, majority-vote and consensus variants. Confident learning prunes instances using out-of-sample predicted probabilities, per-class thresholds and an estimate of the number of label errors \citep{northcutt2021confident}. Our rule is the single-algorithm filter with logistic regression and a plain argmax disagreement. We add a checkable measurement of its precision, recall, accuracy gain and clean-label cost; we do not claim a new filtering idea. We do not compare against confident learning.
